\chapter{संकुलों की अभिक्रिया क्रियाविधियाँ}\label{ch:b3:complex-mechanisms}

ऑक्सीजन-18 से अंकित जल में किसी क्रोमियम(III) लवण को घोलिए, और अंकित जल को
धातु से बँधे जल अणुओं का स्थान लेने में कई दिन लग जाते हैं; कॉपर(II) लवण के साथ
विनिमय मिलाने भर के समय में पूरा हो जाता है। अभिक्रिया वही है, एक जल अणु किसी धातु आयन
को छोड़ता है और दूसरा उसका स्थान लेता है, फिर भी इसकी दर एक आयन से दूसरे तक
अत्यधिक बदलती है, और $d$ इलेक्ट्रॉनों की संख्या इस परिवर्तन का अधिकांश भाग बता देती है।
लिगैंडों के प्रतिस्थापित होने का ढंग तय करता है कि प्लैटिनम कैंसररोधी औषधें कैसे बनती हैं और कैसे
कार्य करती हैं; धातु आयनों के बीच इलेक्ट्रॉनों के चलने का ढंग श्वसन की और हर बैटरी की
गति तय करता है। यह अध्याय संकुलों की अभिक्रियाओं के दो परिवारों का विवेचन करता है:
प्रतिस्थापन और इलेक्ट्रॉन स्थानांतरण, दूसरे का रुडॉल्फ मार्कस के सिद्धांत के साथ।

\begin{recall}
द्वितीय वर्ष के खंड ने संकुलों, लिगैंडों की दंतुरता, सेतु
लिगैंडों, $fac$ और $mer$ समावयवियों, मूल चरण के रूप में लिगैंड विनिमय, और
उच्च- तथा निम्न-चक्रण विन्यासों सहित क्रिस्टल-क्षेत्र स्थायीकरण ऊर्जाओं का वर्णन किया;
प्रथम वर्ष के खंड ने क्रियाविधियों का, उनके स्थायी-अवस्था और दर-निर्धारक-चरण
सन्निकटनों के साथ। \Cref{ch:b3:rate-theories} ने \omterm{thm:b3:rate-theories:eyring}{आयरिंग समीकरण} और
$\Delta^{\ddagger}S^\circ$ का अर्थ दिया, \cref{ch:b3:electrode-kinetics} ने इलेक्ट्रोड पर
इलेक्ट्रॉन स्थानांतरण, और \cref{ch:b3:complex-spectra-magnetism} ने \omterm{def:b3:complex-spectra-magnetism:ligand-field-term}{लिगैंड-क्षेत्र पद} और
यान--टेलर प्रभाव।
\end{recall}

\section{विनिमयशीलता और अक्रियता}

\begin{definition}[विनिमयशील और अक्रिय संकुल]\label{def:b3:complex-mechanisms:labile}
कोई संकुल \emph{विनिमयशील}\index{विनिमयशील संकुल} है यदि उसके लिगैंड
शीघ्र प्रतिस्थापित होते हैं, और \emph{गतिकतः अक्रिय संकुल}\index{गतिकतः अक्रिय संकुल}
यदि वे धीरे प्रतिस्थापित होते हैं; टाउबे के अनुसार, किसी संकुल को विनिमयशील तब कहा जाता है जब
\qty{0.1}{mol/L} और \qty{25}{\celsius} पर किसी लिगैंड के साथ उसकी अभिक्रियाएँ लगभग एक
मिनट में पूरी हो जाती हैं। विनिमयशीलता एक बलगतिक गुणधर्म है, संकुल के ऊष्मागतिक स्थायित्व से
असंबद्ध।
\end{definition}

जल विनिमय, $\mathrm{[M(H_2O)_6]^{n+} + H_2O^{\ast} \to [M(H_2O)_5(H_2O^{\ast})]^{n+} + H_2O}$, सबसे सरल
प्रतिस्थापन है और एक पैमाना तय करता है। मुख्य समूहों के आयन तब तेज़ी से विनिमय करते हैं जब वे
बड़े और कम आवेशित हों। संक्रमण-धातु आयनों में $d^3$
(\ce{Cr^{3+}}) या निम्न-चक्रण $d^6$ (\ce{Co^{3+}}, \ce{Rh^{3+}}, \ce{Ir^{3+}}) विन्यास वाले अक्रिय हैं, और
अधिकांश द्वितीय और तृतीय पंक्ति की धातुएँ भी, जबकि $d^9$ (\ce{Cu^{2+}}) और उच्च-चक्रण $d^4$ (\ce{Cr^{2+}}),
जो यान--टेलर प्रभाव से विकृत हैं, अत्यंत \omterm{def:b3:complex-mechanisms:labile}{विनिमयशील} हैं।

\begin{proposition}[लिगैंड-क्षेत्र सक्रियण ऊर्जा]\label{prop:b3:complex-mechanisms:lfae}
यदि कोई प्रतिस्थापन पंच-उपसहसंयोजी वर्ग-पिरामिडीय संक्रमण
अवस्था से होकर जाता है, तो वहाँ पहुँचने में लिगैंड-क्षेत्र स्थायीकरण की हानि (लिगैंड-क्षेत्र
सक्रियण ऊर्जा), सरलतम $\sigma$-मात्र मॉडल में, निम्न-चक्रण $d^6$
($4\,Dq$) के लिए सबसे बड़ी, फिर $d^3$ और $d^8$ ($2\,Dq$) के लिए, $d^0$, उच्च-चक्रण $d^5$ और $d^{10}$ के लिए शून्य, और
उच्च-चक्रण $d^4$ और $d^9$ के लिए ऋणात्मक होती है।
\end{proposition}

\begin{proof}
प्रत्येक लिगैंड को इस प्रकार प्रतिरूपित कीजिए कि वह किसी $d$ कक्षक को $e_\sigma$ गुणा उस दिशा में कक्षक के
कोणीय आयाम के वर्ग से ऊपर उठाता है, जो पालि अक्ष पर 1 पर प्रसामान्यीकृत है: एक अक्षीय
लिगैंड $d_{z^2}$ को 1 देता है, एक विषुवतीय लिगैंड $d_{z^2}$ को $\frac14$ और $d_{x^2-y^2}$ को $\frac34$; $t_{2g}$ कक्षक
लिगैंडों के बीच इंगित करते हैं और उन्हें कुछ नहीं मिलता। अष्टफलक $d_{z^2}$ और $d_{x^2-y^2}$ प्रत्येक को $3e_\sigma$ देता है
($\Delta_{\mathrm o} = 3e_\sigma = 10Dq$); वर्ग पिरामिड, एक अक्षीय लिगैंड के बिना, $d_{z^2}$ को $2e_\sigma$ और
$d_{x^2-y^2}$ को $3e_\sigma$ देता है। $n$ इलेक्ट्रॉनों का स्थायीकरण उनकी ऊर्जा ऋण $n/5$ गुणा
पाँचों स्तरों का योग (भारकेंद्र) है। निम्न-चक्रण $d^6$ ($t_{2g}^6$) के लिए: अष्टफलक में $0 - \frac65 \times 6e_\sigma$,
पिरामिड में $0 - \frac65 \times 5e_\sigma$: $1.2e_\sigma = 4Dq$ की हानि। $d^3$ के लिए: $\frac35e_\sigma =
2Dq$; $d^8$ ($t_{2g}^6e_g^2$) के लिए: $(5 - \frac85 \times 5) - (6 - \frac85 \times 6) = 0.6e_\sigma$। $d^9$ के लिए: $(7 - 9) - (9 -
10.8) = -0.2e_\sigma$।
\end{proof}

\begin{omfigure}
\begin{tikzpicture}
\begin{axis}[width=0.72\linewidth, height=5.4cm, xmin=-0.6, xmax=10.6, ymin=-1.2, ymax=4.6,
  axis x line=bottom, axis y line=left, ycomb, xtick={0,1,2,3,4,5,6,7,8,9,10},
  xlabel={$n$ ($d^n$)}, ylabel={सक्रियण ऊर्जा / $Dq$}, tick label style={font=\small},
  label style={font=\small}]
\draw[black!40] (axis cs:-0.6,0) -- (axis cs:10.6,0);
\addplot[omDef, line width=4pt, xshift=-2.5pt] table[x=n, y=hs] {figdata/out/bachelor-3-ligand-field-activation.dat};
\addplot[omThm, line width=4pt, xshift=2.5pt, unbounded coords=discard] table[x=n, y=ls] {figdata/out/bachelor-3-ligand-field-activation.dat};
\end{axis}
\end{tikzpicture}

{\small वर्ग पिरामिड से होकर जाने वाले वियोजी पथ के लिए लिगैंड-क्षेत्र सक्रियण ऊर्जाएँ,
प्रस्ताव के $\sigma$-मात्र मॉडल में (नीला: उच्च चक्रण;
लाल: निम्न चक्रण)। सबसे बड़े मान,
निम्न-चक्रण $d^6$, $d^3$ और $d^8$ के लिए, अक्रिय या मंद आयनों (\ce{Co^{3+}}, \ce{Cr^{3+}}, \ce{Ni^{2+}}) से मेल खाते हैं;
ऋणात्मक मान, $d^4$ और $d^9$ के लिए, अत्यंत \omterm{def:b3:complex-mechanisms:labile}{विनिमयशील} यान--टेलर आयनों से।}
\end{omfigure}

\begin{definition}[सक्रियण आयतन]\label{def:b3:complex-mechanisms:activation-volume}
किसी मूल चरण का \emph{सक्रियण आयतन}\index{सक्रियण आयतन} $\Delta V^{\ddagger}$ \omterm{def:b3:rate-theories:activated-complex}{सक्रियित
संकुल} के आंशिक मोलर आयतन और अभिकारकों के आंशिक मोलर आयतनों का अंतर है।
\end{definition}

\begin{proposition}[दर स्थिरांक की दाब-निर्भरता]\label{prop:b3:complex-mechanisms:pressure}
नियत ताप पर, $\displaystyle\Bigl(\frac{\partial\ln k}{\partial p}\Bigr)_T = -\frac{\Delta V^{\ddagger}}{RT}$।
\end{proposition}

\begin{proof}
\omterm{thm:b3:rate-theories:eyring}{आयरिंग समीकरण} से $\ln k = \text{स्थिरांक} - \Delta^{\ddagger}G/RT$, और $(\partial\Delta^{\ddagger}G/\partial p)_T = \Delta V^{\ddagger}$, जैसा
किसी भी गिब्स ऊर्जा के लिए होता है।
\end{proof}

जिस चरण में कोई लिगैंड निकलता है उसके लिए $\Delta V^{\ddagger} > 0$; जिसमें कोई लिगैंड प्रवेश करता है उसके लिए $\Delta V^{\ddagger}
< 0$। कुछ सौ मेगापास्कल पर मापा गया \omterm{def:b3:complex-mechanisms:activation-volume}{सक्रियण आयतन} किसी प्रतिस्थापन क्रियाविधि का
सबसे प्रत्यक्ष निदान है, जो \omterm{def:b3:rate-theories:activation}{सक्रियण एन्ट्रॉपी} की तुलना में विलायकन से कम
धुँधला होता है।

\section{प्रतिस्थापन क्रियाविधियाँ}

\begin{definition}[वियोजी, सहयोजी और विनिमय क्रियाविधियाँ]\label{def:b3:complex-mechanisms:mechanisms}
\emph{वियोजी क्रियाविधि}\index{वियोजी क्रियाविधि} (D) में निर्गामी
लिगैंड पहले निकलता है, जिससे निम्नतर उपसहसंयोजन संख्या का एक मध्यवर्ती बनता है;
\emph{सहयोजी क्रियाविधि}\index{सहयोजी क्रियाविधि} (A) में प्रवेशी लिगैंड
पहले बँधता है, जिससे उच्चतर उपसहसंयोजन संख्या का एक मध्यवर्ती बनता है;
\emph{विनिमय क्रियाविधि}\index{विनिमय क्रियाविधि} (I) में दोनों लिगैंड
एक ही चरण में विनिमय करते हैं, आबंध-विखंडन ($I_d$) या आबंध-निर्माण ($I_a$)
लक्षण के साथ।
\end{definition}

\begin{proposition}[D और A क्रियाविधियों के दर नियम]\label{prop:b3:complex-mechanisms:d-a-laws}
$\mathrm{ML_5X + Y \to ML_5Y + X}$ के लिए: एक D क्रियाविधि ($\mathrm{ML_5X} \rightleftharpoons \mathrm{ML_5} + \mathrm X$, $k_1$, $k_{-1}$; $\mathrm{ML_5} + \mathrm Y
\to \mathrm{ML_5Y}$, $k_2$) देती है
\[
  v = \frac{k_1k_2[\mathrm{ML_5X}][\mathrm Y]}{k_{-1}[\mathrm X] + k_2[\mathrm Y]},
\]
जो संकुल में प्रथम कोटि का है और उच्च $[\mathrm Y]$ पर $[\mathrm Y]$ से स्वतंत्र; स्थायी अवस्था में
एक सप्त-उपसहसंयोजी मध्यवर्ती से होकर जाने वाली A क्रियाविधि $v =
k[\mathrm{ML_5X}][\mathrm Y]$ देती है, प्रत्येक में प्रथम कोटि।
\end{proposition}

\begin{proof}
D: $\mathrm{ML_5}$ के लिए स्थायी अवस्था, $k_1[\mathrm{ML_5X}] = (k_{-1}[\mathrm X] + k_2[\mathrm Y])[\mathrm{ML_5}]$, और $v = k_2[\mathrm{ML_5}][\mathrm Y]$।
उच्च $[\mathrm Y]$ पर $v \to k_1[\mathrm{ML_5X}]$। A: $\mathrm{ML_5X} + \mathrm Y \rightleftharpoons \mathrm{ML_5XY}$ ($k_a$, $k_{-a}$), $\mathrm{ML_5XY} \to
\mathrm{ML_5Y} + \mathrm X$ ($k_b$); स्थायी अवस्था $v = \frac{k_ak_b}{k_{-a} + k_b}[\mathrm{ML_5X}][\mathrm Y]$ देती है।
\end{proof}

जल में प्रवेशी लिगैंड प्रायः विलायक की तुलना में कहीं कम सांद्रता पर उपस्थित होता है,
और वास्तविक पहला चरण संकुल और
लिगैंड की भेंट है।

\begin{definition}[आइगन--विल्किन्स क्रियाविधि]\label{def:b3:complex-mechanisms:eigen-wilkins}
किसी ऐक्वा संकुल के प्रतिस्थापन की \emph{आइगन--विल्किन्स क्रियाविधि}\index{आइगन--विल्किन्स क्रियाविधि}
एक तीव्र पूर्व-साम्य है, जो एक \emph{बाह्य-मंडल संकुल}\index{बाह्य-मंडल संकुल} बनाता है,
जिसमें प्रवेशी लिगैंड संकुल के पास, उसके उपसहसंयोजन मंडल के बाहर, बैठता है, और उसके बाद
एक जल अणु का उस लिगैंड से दर-निर्धारक विनिमय होता है।
\end{definition}

\begin{theorem}[आइगन--विल्किन्स दर नियम]\label{thm:b3:complex-mechanisms:eigen-wilkins}
बाह्य-मंडल साम्य स्थिरांक $K_{\mathrm{os}}$ और विनिमय दर स्थिरांक
$k_i$ के साथ, आधिक्य $[\mathrm Y]$ पर प्रेक्षित छद्म-प्रथम-कोटि दर स्थिरांक है
\[
  k_{\mathrm{obs}} = \frac{K_{\mathrm{os}}k_i[\mathrm Y]}{1 + K_{\mathrm{os}}[\mathrm Y]} .
\]
\end{theorem}

\begin{proof}
संकुल मुक्त और बाह्य-मंडल रूपों के बीच वितरित होता है, $[\mathrm{M\cdot Y}] = K_{\mathrm{os}}[\mathrm M][\mathrm Y]$,
कुल $[\mathrm M]_{\mathrm{tot}} = [\mathrm M](1 + K_{\mathrm{os}}[\mathrm Y])$ के साथ; दर $k_i[\mathrm{M\cdot Y}] = k_iK_{\mathrm{os}}[\mathrm Y][\mathrm M]_{\mathrm{tot}}/(1 + K_{\mathrm{os}}[\mathrm Y])$ है।
\end{proof}

निम्न $[\mathrm Y]$ पर अभिक्रिया द्वितीय कोटि की होती है, $k = K_{\mathrm{os}}k_i$ के साथ; $K_{\mathrm{os}}$, जिसका आकलन
आवेशों और निकटतम पहुँच की दूरी से होता है, एक ही आवेश के सभी लिगैंडों के लिए
लगभग समान होता है, अतः $k_i$, जल-विनिमय चरण, दर को
नियंत्रित करता है: किसी ऐक्वा आयन पर प्रतिस्थापन प्रवेशी लिगैंड से लगभग स्वतंत्र होता है,
जो वियोजी विनिमय का चिह्न है।

\begin{definition}[संयुग्मी-क्षारक क्रियाविधि]\label{def:b3:complex-mechanisms:conjugate-base}
किसी ऐम्मीन संकुल के क्षारीय जल-अपघटन, $\mathrm{[Co(NH_3)_5X]^{2+} + OH^- \to [Co(NH_3)_5(OH)]^{2+} + X^-}$, की
\emph{संयुग्मी-क्षारक क्रियाविधि}\index{संयुग्मी-क्षारक क्रियाविधि}
एक तीव्र पूर्व-साम्य में किसी ऐम्मीन लिगैंड का विप्रोटॉनन है, और उसके बाद
ऐमिडो संयुग्मी क्षारक से $\mathrm X^-$ का वियोजी निष्कासन, जिसका $\mathrm{NH_2^-}$ लिगैंड
संकुल को \omterm{def:b3:complex-mechanisms:labile}{विनिमयशील} बना देता है।
\end{definition}

\begin{proposition}[क्षारीय जल-अपघटन का दर नियम]\label{prop:b3:complex-mechanisms:conjugate-base}
विप्रोटॉनन स्थिरांक $K$ और संयुग्मी क्षारक के वियोजन दर स्थिरांक $k$ के
साथ, $v = \dfrac{kK[\text{संकुल}]_{\mathrm{tot}}[\ce{OH-}]}{1 + K[\ce{OH-}]}$, निम्न $[\ce{OH-}]$ पर द्वितीय कोटि।
\end{proposition}

\begin{proof}
आइगन--विल्किन्स नियम की तरह: संकुल अपने अम्ल और
संयुग्मी-क्षारक रूपों के बीच $1 : K[\ce{OH-}]$ के अनुपात में वितरित होता है, और केवल दूसरा अभिक्रिया करता है।
\end{proof}

केवल दर नियम इस क्रियाविधि को \ce{OH-} के प्रत्यक्ष आक्रमण से
अलग नहीं करता; प्रमाण यह है कि क्षारीय जल-अपघटन के लिए N--H प्रोटॉन आवश्यक है (जिन संकुलों में
यह नहीं होता वे इसे नहीं दिखाते) और ऐम्मीन प्रोटॉन जल-अपघटन की तुलना में विलायक से कहीं
तेज़ी से विनिमय करते हैं।

\begin{method}[प्रतिस्थापन क्रियाविधि का निदान]\label{met:b3:complex-mechanisms:diagnose}
\begin{enumerate}
  \item प्रवेशी-समूह निर्भरता: Y से स्वतंत्र (या संतृप्त होती) दर,
        वियोजी; $[\mathrm Y]$ के समानुपाती और Y की प्रकृति के प्रति संवेदी दर,
        सहयोजी।
  \item \omterm{def:b3:complex-mechanisms:activation-volume}{सक्रियण आयतन}: धनात्मक, वियोजी; ऋणात्मक, सहयोजी।
  \item \omterm{def:b3:rate-theories:activation}{सक्रियण एन्ट्रॉपी}: D के लिए धनात्मक, A के लिए ऋणात्मक (सावधानी के साथ, क्योंकि
        विलायकन का योगदान होता है)।
  \item निर्गामी-समूह निर्भरता: M--X आबंध की प्रबलता का अनुसरण करने वाली दर
        संक्रमण अवस्था में आबंध-विखंडन की ओर संकेत करती है।
\end{enumerate}
\end{method}

\begin{omfigure}
\begin{tikzpicture}
\begin{axis}[name=pd, omprofile, width=0.46\linewidth, height=4.6cm, xmin=0, xmax=10, ymin=0, ymax=10,
  xlabel={अभिक्रिया निर्देशांक}, ylabel={ऊर्जा}, title={\small D: $\mathrm{ML_5X} \to \mathrm{ML_5} + \mathrm X \to \mathrm{ML_5Y}$}]
\addplot[omDef, very thick, smooth] coordinates {(0,2) (1.5,2.1) (3,7.5) (4.2,5.6) (5.0,5.4) (5.8,5.6) (7,7.0) (8.5,1.6) (10,1.5)};
\node[font=\scriptsize, anchor=north, align=center] at (axis cs:5,5.2) {ML$_5$ (CN 5)};
\end{axis}
\begin{axis}[at={(pd.east)}, anchor=west, xshift=1.0cm, omprofile, width=0.46\linewidth, height=4.6cm,
  xmin=0, xmax=10, ymin=0, ymax=10, xlabel={अभिक्रिया निर्देशांक}, ylabel={ऊर्जा},
  title={\small A: $\mathrm{ML_5X} + \mathrm Y \to \mathrm{ML_5XY} \to \mathrm{ML_5Y}$}]
\addplot[omThm, very thick, smooth] coordinates {(0,2) (1.5,2.1) (3,6.5) (4.2,4.4) (5.0,4.2) (5.8,4.4) (7,7.2) (8.5,1.6) (10,1.5)};
\node[font=\scriptsize, anchor=north, align=center] at (axis cs:5,4.0) {ML$_5$XY (CN 7)};
\end{axis}
\end{tikzpicture}

{\small एक वियोजी और एक सहयोजी प्रतिस्थापन के ऊर्जा प्रोफ़ाइल
(व्यवस्थात्मक): दोनों एक मध्यवर्ती से होकर जाते हैं, D में निम्नतर उपसहसंयोजन संख्या
(CN) का और A में उच्चतर उपसहसंयोजन संख्या का। \omterm{def:b3:complex-mechanisms:mechanisms}{विनिमय क्रियाविधियों} में
कूप लुप्त हो जाता है और एक ही संक्रमण अवस्था रहती है।}
\end{omfigure}

वर्ग पिरामिड से होकर प्रतिस्थापन करने वाला अष्टफलकीय संकुल सामान्यतः
अपना विन्यास बनाए रखता है ($cis$, $cis$ ही रहता है); जो त्रिकोणीय
द्विपिरामिड से होकर जाता है, जैसा अनेक कोबाल्ट(III) संकुल क्षारीय जल-अपघटन में करते हैं, वह
$cis$ और $trans$ उत्पादों का मिश्रण दे सकता है।

\section{वर्ग-समतलीय प्रतिस्थापन और ट्रांस प्रभाव}

$d^8$ आयनों (\ce{Pt^{2+}}, \ce{Pd^{2+}}, \ce{Au^{3+}}, प्रबल लिगैंडों के साथ \ce{Ni^{2+}}) के वर्ग-समतलीय संकुल
उपसहसंयोजी रूप से असंतृप्त होते हैं: प्रवेशी लिगैंड रिक्त अक्षीय दिशा के अनुदिश
पास आ सकता है। उनके प्रतिस्थापन सहयोजी होते हैं, एक
पंच-उपसहसंयोजी त्रिकोणीय-द्विपिरामिडीय संक्रमण अवस्था से होकर।

\begin{proposition}[द्वि-पद दर नियम]\label{prop:b3:complex-mechanisms:square-planar}
किसी उपसहसंयोजी विलायक S में प्रतिस्थापन $\mathrm{[PtL_3X] + Y \to [PtL_3Y] + X}$, आधिक्य Y पर,
$k_{\mathrm{obs}} = k_1 + k_2[\mathrm Y]$ का पालन करता है।
\end{proposition}

\begin{proof}
दो समांतर सहयोजी पथ: Y का प्रत्यक्ष आक्रमण ($k_2[\mathrm Y]$), और
बड़े आधिक्य में उपस्थित विलायक का आक्रमण (छद्म-प्रथम-कोटि $k_1$), जो $[\mathrm{PtL_3S}]$ देता है, जिसे
फिर Y शीघ्र रूपांतरित कर देता है। समांतर प्रथम-कोटि पथ जुड़ जाते हैं।
\end{proof}

\begin{definition}[ट्रांस प्रभाव, ट्रांस अनुप्रभाव]\label{def:b3:complex-mechanisms:trans-effect}
\emph{ट्रांस प्रभाव}\index{ट्रांस प्रभाव} किसी लिगैंड के प्रतिस्थापन का उसके
ट्रांस स्थित लिगैंड द्वारा त्वरण है: संक्रमण अवस्था पर एक बलगतिक प्रभाव।
\emph{ट्रांस अनुप्रभाव}\index{ट्रांस अनुप्रभाव} मूल अवस्था में किसी लिगैंड के
ट्रांस स्थित आबंध का दुर्बल होना है, जो आबंध लंबाइयों और कंपन
आवृत्तियों में दिखता है: एक ऊष्मागतिक प्रभाव।
\end{definition}

\omterm{def:b3:complex-mechanisms:trans-effect}{ट्रांस प्रभाव} इस क्रम का पालन करता है: $\ce{CO}$, $\ce{CN-}$, $\ce{C2H4} > \ce{PR3}$, $\ce{H-} > \ce{CH3-} > \ce{I-}$, $\ce{SCN-} >
\ce{Br-} > \ce{Cl-} > \ce{NH3}$, $\ce{H2O}$। प्रबल $\sigma$ दाता ट्रांस आबंध को दुर्बल करते हैं (\omterm{def:b3:complex-mechanisms:trans-effect}{ट्रांस अनुप्रभाव});
प्रबल $\pi$ ग्राही धातु से \omterm{def:b3:computational-chemistry:dft}{इलेक्ट्रॉन घनत्व} लेकर पंच-उपसहसंयोजी संक्रमण अवस्था को
स्थायी बनाते हैं। दोनों दर बढ़ाते हैं।

\begin{method}[प्लैटिनम(II) समावयवियों के संश्लेषण की योजना]\label{met:b3:complex-mechanisms:pt-synthesis}
\begin{enumerate}
  \item प्रत्येक चरण में वह लिगैंड प्रतिस्थापित होता है जो उपस्थित सर्वाधिक \omterm{def:b3:complex-mechanisms:trans-effect}{ट्रांस प्रभाव} वाले
        लिगैंड के ट्रांस स्थित है।
  \item योगों का क्रम ऐसा रखिए कि यह नियम वांछित समावयवी तक ले जाए।
\end{enumerate}
\end{method}

\begin{omfigure}
\setchemfig{atom sep=2.2em}
\schemestart
\chemfig{Cl-Pt(-[2]Cl)(-[6]Cl)-Cl}
\arrow{->[\ce{NH3}]}
\chemfig{Cl-Pt(-[2]Cl)(-[6]Cl)-NH_3}
\arrow{->[\ce{NH3}]}
\chemfig{Cl-Pt(-[2]NH_3)(-[6]Cl)-NH_3}
\schemestop

\bigskip
\schemestart
\chemfig{H_3N-Pt(-[2]NH_3)(-[6]NH_3)-NH_3}
\arrow{->[\ce{Cl-}]}
\chemfig{H_3N-Pt(-[2]NH_3)(-[6]NH_3)-Cl}
\arrow{->[\ce{Cl-}]}
\chemfig{Cl-Pt(-[2]NH_3)(-[6]NH_3)-Cl}
\schemestop

{\small \omterm{def:b3:complex-mechanisms:trans-effect}{ट्रांस प्रभाव} कार्यरत (आवेश छोड़ दिए गए हैं)। ऊपर: $\ce{[PtCl4]^{2-}}$ से, दूसरा
अमोनिया क्लोराइड के ट्रांस स्थित एक क्लोराइड को प्रतिस्थापित करता है (Cl$^-$ का \omterm{def:b3:complex-mechanisms:trans-effect}{ट्रांस प्रभाव} अधिक है),
जिससे $cis$ समावयवी, सिसप्लैटिन, बनता है। नीचे: $\ce{[Pt(NH3)4]^{2+}}$ से, दूसरा क्लोराइड
पहले क्लोराइड के ट्रांस स्थित अमोनिया को प्रतिस्थापित करता है, जिससे $trans$ समावयवी बनता है।}
\end{omfigure}

\section{इलेक्ट्रॉन स्थानांतरण}

\begin{definition}[बाह्य- और आंतर-मंडल इलेक्ट्रॉन स्थानांतरण]\label{def:b3:complex-mechanisms:electron-transfer}
\emph{बाह्य-मंडल इलेक्ट्रॉन स्थानांतरण}\index{बाह्य-मंडल इलेक्ट्रॉन स्थानांतरण} में
इलेक्ट्रॉन ऐसे दो संकुलों के बीच जाता है जिनके उपसहसंयोजन मंडल अक्षुण्ण रहते हैं;
\emph{आंतर-मंडल इलेक्ट्रॉन स्थानांतरण}\index{आंतर-मंडल इलेक्ट्रॉन स्थानांतरण} में यह
एक क्षणिक द्विनाभिकीय संकुल में दोनों धातुओं को जोड़ने वाले सेतु लिगैंड से होकर
जाता है। \emph{स्व-विनिमय अभिक्रिया}\index{स्व-विनिमय अभिक्रिया} एक ही युग्म की दो
ऑक्सीकरण अवस्थाओं के बीच इलेक्ट्रॉन स्थानांतरण है, बिना किसी
नेट रासायनिक परिवर्तन के, जैसे $\ce{[Fe(H2O)6]^{3+} + [Fe(H2O)6]^{2+}}$।
\end{definition}

हेनरी टाउबे ने 1953 में आंतर-मंडल पथ को एक ऐसी अभिक्रिया से दिखाया जो इस प्रकार चुनी गई थी कि
उत्पाद स्वयं अपनी कहानी बता दे। कोबाल्ट(III) अक्रिय है, कोबाल्ट(II) \omterm{def:b3:complex-mechanisms:labile}{विनिमयशील}; क्रोमियम(II)
\omterm{def:b3:complex-mechanisms:labile}{विनिमयशील} है, क्रोमियम(III) अक्रिय। जब $\ce{[Co(NH3)5Cl]^{2+}}$ को अम्ल में $\ce{[Cr(H2O)6]^{2+}}$ द्वारा
अपचयित किया जाता है, तो बना सारा क्रोमियम(III) क्लोराइड को धारण करता है, $\ce{[Cr(H2O)5Cl]^{2+}}$ के रूप में, मुक्त
क्लोराइड की उपस्थिति में भी: क्लोराइड ने अवश्य दोनों धातुओं को जोड़ा होगा,
इलेक्ट्रॉन के गुज़रते ही अक्रिय क्रोमियम(III) से बंध गया होगा, और
\omterm{def:b3:complex-mechanisms:labile}{विनिमयशील} कोबाल्ट(II) द्वारा मुक्त कर दिया गया होगा।

\begin{omfigure}
\begin{tikzpicture}[font=\small]
  \begin{scope}[scale=0.8]
    \pic[transform shape] (a) at (0,0) {omoctahedron};
    \pic[transform shape] (b) at (2.0,0) {omoctahedron};
  \end{scope}
  \node[anchor=north east] at (0.3,-1.0) {$\mathrm{Co^{III}(NH_3)_5}$};
  \node[anchor=north west] at (1.3,-1.0) {$\mathrm{Cr^{II}(H_2O)_5}$};
  \node[omThm, anchor=south] at (0.8,0.15) {Cl};
  \draw[->, thick, omDef] (1.6,1.75) to[bend right=40] node[midway, above] {e$^-$} (0,1.75);
\end{tikzpicture}

{\small टाउबे का सेतुबद्ध मध्यवर्ती (व्यवस्थात्मक): क्लोराइड, जो अब भी
कोबाल्ट(III) से बँधा है, एक जल अणु को प्रतिस्थापित करके क्रोमियम(II) के उपसहसंयोजन मंडल में
प्रवेश करता है; इलेक्ट्रॉन सेतु से होकर पार जाता है, और क्लोराइड
अब अक्रिय क्रोमियम(III) पर रह जाता है।}
\end{omfigure}

एक इलेक्ट्रॉन लगभग $10^{-16}$ s में छलाँग लगाता है, जो नाभिकों की गति से कहीं तेज़ है:
\omterm{thm:b3:electronic-spectroscopy:franck-condon}{फ़्रैंक--कॉन्डन सिद्धांत} (\cref{ch:b3:electronic-spectroscopy}) के अनुसार स्थानांतरण
केवल ऐसे नाभिकीय विन्यास पर हो सकता है जहाँ अभिकारकों और उत्पादों की
ऊर्जा समान हो। आयरन स्व-विनिमय के लिए \ce{Fe^{III}}--O आबंध
\ce{Fe^{II}}--O आबंधों से छोटे हैं; इलेक्ट्रॉन के चलने से पहले दोनों संकुलों को ऊर्जा की कीमत पर
एक साझा मध्यवर्ती ज्यामिति तक विकृत होना पड़ता है।

\section{मार्कस सिद्धांत}

\begin{definition}[पुनर्गठन ऊर्जा]\label{def:b3:complex-mechanisms:reorganisation}
किसी इलेक्ट्रॉन स्थानांतरण की \emph{पुनर्गठन ऊर्जा}\index{पुनर्गठन ऊर्जा} $\lambda$ वह
ऊर्जा है जो अभिकारकों को, इलेक्ट्रॉन का स्थानांतरण किए बिना, उत्पादों के
नाभिकीय विन्यास (संकुलों की आबंध लंबाइयाँ और आसपास के विलायक के अभिविन्यास)
तक लाने के लिए आवश्यक है। यह आबंधों से आने वाले एक आंतरिक भाग और विलायक से आने वाले
एक बाह्य भाग का योग है।
\end{definition}

\begin{theorem}[मार्कस समीकरण]\label{thm:b3:complex-mechanisms:marcus}
यदि अभिकारकों और उत्पादों की गिब्स ऊर्जाएँ किसी अभिक्रिया निर्देशांक $x$ में (अभिकारकों के साम्य पर 0,
उत्पादों के साम्य पर 1) समान वक्रता के परवलय हों, तो
मानक अभिक्रिया गिब्स ऊर्जा $\Delta G^\circ$ और \omterm{def:b3:complex-mechanisms:reorganisation}{पुनर्गठन ऊर्जा} $\lambda$ वाले इलेक्ट्रॉन स्थानांतरण की
\omterm{def:b3:rate-theories:activation}{सक्रियण गिब्स ऊर्जा} है
\[
  \Delta G^{\ddagger} = \frac{(\lambda + \Delta G^\circ)^2}{4\lambda} .
\]
यह \emph{मार्कस समीकरण}\index{मार्कस समीकरण} है।
\end{theorem}

\begin{proof}
$G_R = \lambda x^2$ और $G_P = \lambda(x - 1)^2 + \Delta G^\circ$ लिखिए: वक्रता $G_R(1) = \lambda$ से तय होती है, जो
$\lambda$ की परिभाषा है। स्थानांतरण वहाँ होता है जहाँ वक्र एक-दूसरे को काटते हैं: $\lambda x^2 = \lambda x^2 - 2\lambda x + \lambda +
\Delta G^\circ$, अतः $x^\ast = (\lambda + \Delta G^\circ)/2\lambda$, और $\Delta G^{\ddagger} = G_R(x^\ast) = (\lambda + \Delta G^\circ)^2/4\lambda$।
\end{proof}

\begin{corollary}[प्रतिलोमित क्षेत्र]\label{cor:b3:complex-mechanisms:inverted}
नियत $\lambda$ पर, अभिक्रिया के अधिक ऊर्जाक्षेपी होने के साथ दर तब तक बढ़ती है जब तक
$-\Delta G^\circ = \lambda$, जहाँ $\Delta G^{\ddagger} = 0$; उसके आगे यह फिर घटती है।
\end{corollary}

\begin{proof}
$\Delta G^{\ddagger}$, $\Delta G^\circ$ में एक परवलय है, जिसका निम्निष्ठ, शून्य, $\Delta G^\circ = -\lambda$ पर है; $\Delta G^\circ < -\lambda$ के लिए यह
फिर बढ़ता है।
\end{proof}

\begin{definition}[प्रतिलोमित क्षेत्र]\label{def:b3:complex-mechanisms:inverted}
इलेक्ट्रॉन स्थानांतरण का \emph{प्रतिलोमित क्षेत्र}\index{प्रतिलोमित क्षेत्र}
प्रेरक बलों का वह परास $-\Delta G^\circ > \lambda$ है जिसमें अभिक्रिया के अधिक ऊर्जाक्षेपी होने पर
दर घटती है।
\end{definition}

\begin{omfigure}
\begin{tikzpicture}
\begin{axis}[name=mp, width=0.47\linewidth, height=5.6cm, xmin=-0.6, xmax=2.2, ymin=-160, ymax=120,
  axis lines=left, xlabel={अभिक्रिया निर्देशांक $x$}, ylabel={$G$ / \unit{kJ.mol^{-1}}},
  tick label style={font=\small}, label style={font=\small}, xtick={0,1,2}]
\addplot[black, thick] table[x=x, y=GR] {figdata/out/bachelor-3-marcus-parabolas.dat};
\addplot[omDef!50, thick] table[x=x, y=P0] {figdata/out/bachelor-3-marcus-parabolas.dat};
\addplot[omDef, thick] table[x=x, y=P50] {figdata/out/bachelor-3-marcus-parabolas.dat};
\addplot[omThm, thick] table[x=x, y=P100] {figdata/out/bachelor-3-marcus-parabolas.dat};
\addplot[omThm!60, thick, dashed] table[x=x, y=P150] {figdata/out/bachelor-3-marcus-parabolas.dat};
\node[font=\scriptsize, anchor=north] at (axis cs:0,-2) {R};
\end{axis}
\begin{axis}[at={(mp.east)}, anchor=west, xshift=1.4cm, width=0.47\linewidth, height=5.6cm,
  xmin=0, xmax=250, ymin=0, ymax=11.5, axis lines=left, xlabel={$-\Delta G^\circ$ / \unit{kJ.mol^{-1}}},
  ylabel={$\log(k/\unit{s^{-1}})$}, tick label style={font=\small}, label style={font=\small}]
\addplot[omDef, very thick] table[x=mdG, y=logk] {figdata/out/bachelor-3-marcus-rate.dat};
\draw[black!40, dashed] (axis cs:100,0) -- (axis cs:100,11);
\node[font=\scriptsize, anchor=south east] at (axis cs:95,1) {सामान्य};
\node[font=\scriptsize, anchor=south west] at (axis cs:105,1) {प्रतिलोमित};
\end{axis}
\end{tikzpicture}

{\small $\lambda = \qty{100}{kJ/mol}$ के साथ मार्कस सिद्धांत (मॉडल)। बाएँ: अभिकारक परवलय
(काला) और उत्पाद परवलय, $\Delta G^\circ = 0$ और $-50$ (नीला), $-100 = -\lambda$ (लाल) तथा
\qty{-150}{kJ/mol} (धराशायी, \omterm{def:b3:complex-mechanisms:inverted}{प्रतिलोमित क्षेत्र}) के लिए; उनका प्रतिच्छेदन
बिंदु, संक्रमण अवस्था, $-\Delta G^\circ = \lambda$ पर अभिकारक वक्र की तली तक
गिरता है और उसके आगे फिर उठता है। दाएँ: दर स्थिरांक का लघुगणक
(मॉडल पूर्व-गुणक \qty{e11}{s^{-1}}) $-\Delta G^\circ = \lambda$ पर एक अधिकतम से गुज़रता है।}
\end{omfigure}

\begin{proposition}[बाह्य-मंडल पुनर्गठन ऊर्जा]\label{prop:b3:complex-mechanisms:outer-reorganisation}
त्रिज्याओं $a_1$, $a_2$ वाले दो गोलाकार अभिकारकों के लिए, जो दूरी $d$ पर हैं, अपवर्तनांक
$n$ और आपेक्षिक परावैद्युतांक $\varepsilon_s$ वाले विलायक में, $\lambda$ का विलायक भाग
$\bigl(\frac{1}{2a_1} + \frac{1}{2a_2} - \frac1d\bigr)\bigl(\frac{1}{n^2} - \frac{1}{\varepsilon_s}\bigr)$ के समानुपाती है।
\end{proposition}

\admitted

जल, जो अत्यंत ध्रुवीय है ($\varepsilon_s$ बड़ा) पर जिसका अपवर्तनांक सामान्य है, बड़ा
विलायक पुनर्गठन देता है; अध्रुवीय विलायक छोटा। बड़े अभिकारक
कम पुनर्गठित होते हैं: यही कारण है कि इलेक्ट्रॉन-स्थानांतरण प्रोटीन अपने धातु स्थलों को
प्रोटीन के भीतर दबाए रखते हैं।

\begin{theorem}[मार्कस संकर संबंध]\label{thm:b3:complex-mechanisms:cross-relation}
किसी संकर अभिक्रिया $\mathrm{A_{ox} + B_{red} \to A_{red} + B_{ox}}$ के लिए, जिसका साम्य स्थिरांक $K_{12}$ है और जो ऐसे युग्मों के बीच है
जिनके स्व-विनिमय दर स्थिरांक $k_{11}$ और $k_{22}$ हैं,
\[
  k_{12} = \sqrt{k_{11}k_{22}K_{12}f}, \qquad \ln f = \frac{(\ln K_{12})^2}{4\ln(k_{11}k_{22}/Z^2)},
\]
जहाँ $Z$ संघट्ट आवृत्ति गुणक है; $f \approx 1$, जब $K_{12}$ बहुत बड़ा न हो। यह
\emph{मार्कस संकर संबंध}\index{मार्कस संकर संबंध} है।
\end{theorem}

\begin{proof}[आंशिक उपपत्ति]
मान लीजिए संकर अभिक्रिया की \omterm{def:b3:complex-mechanisms:reorganisation}{पुनर्गठन ऊर्जा} स्व-विनिमयों की \omterm{def:b3:complex-mechanisms:reorganisation}{पुनर्गठन ऊर्जाओं} का
माध्य है, $\lambda_{12} = (\lambda_{11} + \lambda_{22})/2$ (प्रत्येक अभिकारक अपने स्व-विनिमय
पुनर्गठन का आधा योगदान देता है), और सभी दर स्थिरांक $k = Z\eu^{-\Delta G^{\ddagger}/RT}$। तब
$\Delta G_{12}^{\ddagger} = \frac{\lambda_{12}}{4}\bigl(1 + \frac{\Delta G^\circ_{12}}{\lambda_{12}}\bigr)^2 = \frac{\lambda_{12}}{4} + \frac{\Delta G^\circ_{12}}{2} + \frac{(\Delta G^\circ_{12})^2}{4\lambda_{12}}$। पहला पद
$\sqrt{k_{11}k_{22}}$ देता है (क्योंकि $\Delta G^{\ddagger}_{ii} = \lambda_{ii}/4$), दूसरा $\sqrt{K_{12}}$ ($\Delta G^\circ_{12} = -RT\ln K_{12}$ के साथ), तीसरा
गुणक $f$।
\end{proof}

\begin{method}[स्व-विनिमय आँकड़ों से संकर-अभिक्रिया दर]\label{met:b3:complex-mechanisms:marcus-estimate}
\begin{enumerate}
  \item दोनों युग्मों के स्व-विनिमय स्थिरांक $k_{11}$, $k_{22}$ ज्ञात कीजिए।
  \item मानक विभवों से $K_{12}$ परिकलित कीजिए: $\ln K_{12} = nF(E^\circ_{\text{ऑक्सीकारक}} - E^\circ_{\text{अपचायक}})/RT$।
  \item $k_{12} = \sqrt{k_{11}k_{22}K_{12}}$ ($f$ के साथ, यदि $K_{12}$ बहुत बड़ा हो)। मापे गए मान से कुछ गुना के
        भीतर सहमति बाह्य-मंडल क्रियाविधि का संकेत देती है।
\end{enumerate}
\end{method}

मार्कस ने यह सिद्धांत 1956 में प्रकाशित किया; \omterm{def:b3:complex-mechanisms:inverted}{प्रतिलोमित क्षेत्र}, जिस पर पहले संदेह किया गया था,
1984 में गेरहार्ड क्लॉस और जॉन मिलर ने ऐसे अणुओं में प्रेक्षित किया जिनमें दाता और
ग्राही एक दृढ़ अंतरक द्वारा नियत दूरी पर रखे जाते हैं, ताकि विसरण
उसे छिपा न सके।

\begin{inthelab}[ऑक्सीजन-17 NMR द्वारा जल विनिमय]
किसी अनुचुंबकीय आयन से बँधे जल की ऑक्सीजन-17 अनुनाद और स्थूल
जल की अनुनाद उनके बीच विनिमय से चौड़ी हो जाती हैं। \ce{^{17}O} से समृद्ध धातु परक्लोरेट के विलयन में
स्थूल जल संकेत की रेखा-चौड़ाई को ताप के विरुद्ध अभिलिखित करने से
विनिमय दर स्थिरांक मिलता है और, उच्च-दाब प्रोब में,
उसका \omterm{def:b3:complex-mechanisms:activation-volume}{सक्रियण आयतन}।
\end{inthelab}

\begin{safety}{\ghs[0.8cm]{GHS05}\ghs[0.8cm]{GHS06}\ghs[0.8cm]{GHS08}}
% ledger: ghs4:K2PtCl4
पोटैशियम टेट्राक्लोरोप्लैटिनेट(II): निगलने पर विषैला, आँखों को गंभीर
क्षति पहुँचाता है, श्वसन और त्वचा सुग्राहीकारक। प्लैटिनम संकुल
धूम-हुड में दस्ताने पहनकर सँभाले जाते हैं; सिसप्लैटिन और उसके अनुरूप कोशिकाविषी औषधें हैं और
उन्हें उसी रूप में सँभाला जाता है।
\end{safety}

\begin{history}[टाउबे और मार्कस]
हेनरी टाउबे ने 1952 में संकुलों को \omterm{def:b3:complex-mechanisms:labile}{विनिमयशील} या अक्रिय के रूप में वर्गीकृत किया, और अपने
1953 के क्रोमियम--कोबाल्ट प्रयोग से आंतर-मंडल क्रियाविधि स्थापित की; उन्हें
रसायन का 1983 का नोबेल पुरस्कार मिला। रुडॉल्फ मार्कस ने 1956 से अपने नाम वाला
इलेक्ट्रॉन स्थानांतरण सिद्धांत व्युत्पन्न किया, और उन्हें रसायन का 1992 का नोबेल
पुरस्कार मिला।
\end{history}

\section{अभ्यास}

\begin{exercise}[$\star$]\label{exo:b3:complex-mechanisms:1}
कारण सहित \omterm{def:b3:complex-mechanisms:labile}{विनिमयशील} या अक्रिय के रूप में वर्गीकृत कीजिए: $\ce{[Cr(H2O)6]^{3+}}$, $\ce{[Co(NH3)6]^{3+}}$, $\ce{[Cr(H2O)6]^{2+}}$,
$\ce{[Cu(H2O)6]^{2+}}$, $\ce{[Zn(H2O)6]^{2+}}$।
\end{exercise}

\begin{exercise}[$\star$]\label{exo:b3:complex-mechanisms:2}
दिखाइए कि D दर नियम $v = k_1[\mathrm{ML_5X}]$ में बदल जाता है, उच्च $[\mathrm Y]$ पर, और $v = (k_1k_2/k_{-1})
[\mathrm{ML_5X}][\mathrm Y]/[\mathrm X]$ में, जब $k_{-1}[\mathrm X] \gg k_2[\mathrm Y]$।
\end{exercise}

\begin{exercise}[$\star$]\label{exo:b3:complex-mechanisms:3}
किसी D और किसी A प्रतिस्थापन के लिए $\Delta V^{\ddagger}$ और $\Delta^{\ddagger}S^\circ$ के चिह्नों का पूर्वानुमान कीजिए।
\end{exercise}

\begin{exercise}[$\star$]\label{exo:b3:complex-mechanisms:4}
$\ce{[PtCl4]^{2-}}$ के दो \ce{NH3} के साथ, और $\ce{[Pt(NH3)4]^{2+}}$ के दो \ce{Cl-} के साथ उत्पाद बताइए।
\end{exercise}

\begin{exercise}[$\star\star$]\label{exo:b3:complex-mechanisms:5}
किसी ऐक्वा आयन पर एक प्रतिस्थापन के लिए $K_{\mathrm{os}} = \qty{2.0}{L.mol^{-1}}$ और $k_i = \qty{3.0e4}{s^{-1}}$ (अभ्यास के
आँकड़े)। $k_{\mathrm{obs}}$ परिकलित कीजिए, $[\mathrm Y] = 0.10$ और \qty{1.0}{mol/L} पर, और उसकी सीमा।
\end{exercise}

\begin{exercise}[$\star\star$]\label{exo:b3:complex-mechanisms:6}
\qty{298}{K} पर दाब 0.1 से \qty{100}{MPa} तक बढ़ने पर एक दर स्थिरांक दुगुना हो जाता है।
$\Delta V^{\ddagger}$ परिकलित कीजिए और निष्कर्ष निकालिए।
\end{exercise}

\begin{exercise}[$\star\star$]\label{exo:b3:complex-mechanisms:7}
$\ce{[PtCl3(PEt3)]-}$ में \ce{PEt3} के ट्रांस स्थित Pt--Cl आबंध \qty{2.42}{\angstrom} है और अन्य दो \qty{2.32}{\angstrom}
(अभ्यास के आँकड़े)। व्याख्या कीजिए, और पूर्वानुमान कीजिए कि कौन-सा क्लोराइड पहले प्रतिस्थापित होगा।
\end{exercise}

\begin{exercise}[$\star\star$]\label{exo:b3:complex-mechanisms:8}
उन प्रायोगिक प्रमाणों की सूची बनाइए जो किसी इलेक्ट्रॉन स्थानांतरण को आंतर-मंडल
सिद्ध करेंगे।
\end{exercise}

\begin{exercise}[$\star\star$]\label{exo:b3:complex-mechanisms:9}
$\lambda = \qty{80}{kJ/mol}$ के लिए, $\Delta G^{\ddagger}$ परिकलित कीजिए, $\Delta G^\circ = 0$, $-20$, $-80$ और \qty{-120}{kJ/mol} के लिए।
\end{exercise}

\begin{exercise}[$\star\star\star$]\label{exo:b3:complex-mechanisms:10}
दो युग्मों के स्व-विनिमय स्थिरांक $k_{11} = 4.0$ और $k_{22} = \qty{2.0e5}{L.mol^{-1}.s^{-1}}$ हैं, और
उनकी संकर अभिक्रिया के लिए $K_{12} = \num{1.0e4}$ (अभ्यास के आँकड़े)। $k_{12}$ का आकलन कीजिए ($f = 1$)।
\end{exercise}

\begin{exercise}[$\star\star\star$]\label{exo:b3:complex-mechanisms:11}
$\sigma$-मात्र मॉडल का उपयोग करके उच्च-चक्रण $d^5$, $d^7$ और $d^8$ आयनों की लिगैंड-क्षेत्र सक्रियण ऊर्जाएँ
परिकलित कीजिए और \ce{Mn^{2+}}, \ce{Co^{2+}} और \ce{Ni^{2+}} को अपेक्षित
जल-विनिमय दर के क्रम में रखिए।
\end{exercise}

\begin{exercise}[$\star\star\star$]\label{exo:b3:complex-mechanisms:12}
विलयन में विसरण द्वारा पास आने वाले आयनों के बीच की अभिक्रियाओं के लिए \omterm{def:b3:complex-mechanisms:inverted}{प्रतिलोमित क्षेत्र} का प्रेक्षण
कठिन क्यों है, और दृढ़ दाता--ग्राही अणुओं ने उसे कैसे प्रकट
किया?
\end{exercise}

\section{समस्या: सिसप्लैटिन बनाना, ट्रांसप्लैटिन नहीं}

\begin{problem}[{सप्ताहांत समस्या --- सिसप्लैटिन बनाना, ट्रांसप्लैटिन नहीं: वर्ग-समतलीय
प्लैटिनम(II) आयन, ट्रांस प्रभाव से नियोजित एक संश्लेषण, औषध का
ऐक्वाकरण, और कोशिकाओं के भीतर उसके सक्रिय होने का कारण}]
\label{pb:b3:complex-mechanisms:1}
% ledger: ghs4:K2PtCl4
सिसप्लैटिन, $cis$-$\ce{[PtCl2(NH3)2]}$, $\ce{K2[PtCl4]}$ से बनाया जाता है। समस्या के आँकड़े:
\qty{37}{\celsius} पर प्रथम ऐक्वाकरण, $\ce{[PtCl2(NH3)2] + H2O -> [PtCl(H2O)(NH3)2]+ + Cl-}$, का
दर स्थिरांक $k = \qty{9.0e-5}{s^{-1}}$ और साम्य स्थिरांक $K = \qty{3.0e-3}{mol/L}$ है;
क्लोराइड सांद्रता रक्त प्लाज़्मा में लगभग \qty{100}{mM} और कोशिकाओं के भीतर \qty{4}{mM} है।

\medskip\noindent\textbf{भाग I --- प्लैटिनम(II)।}
\begin{enumerate}
  \item \ce{Pt^{2+}} की $d$-इलेक्ट्रॉन संख्या बताइए।
  \item इसके संकुल वर्ग-समतलीय और प्रतिचुंबकीय क्यों होते हैं?
  \item उनके प्रतिस्थापन सहयोजी क्यों होते हैं?
  \item द्वि-पद दर नियम लिखिए और उसके दोनों पदों की व्याख्या कीजिए।
  \item एक दिन के समय-पैमाने पर औषध \omterm{def:b3:complex-mechanisms:labile}{विनिमयशील} है या अक्रिय?
\end{enumerate}

\medskip\noindent\textbf{भाग II --- संश्लेषण।}
\begin{enumerate}[resume]
  \item $\ce{K2[PtCl4]}$ दो \ce{NH3} के साथ क्या देता है? क्यों?
  \item वह प्रत्यक्ष मार्ग अशुद्ध उत्पाद देता है। धारा का मार्ग पहले
        $\ce{[PtCl4]^{2-}}$ को आधिक्य KI के साथ $\ce{[PtI4]^{2-}}$ में बदलता है। आयोडाइड उपयोगी क्यों है?
  \item \ce{NH3} मिलाने से $cis$-$\ce{[PtI2(NH3)2]}$ मिलता है। त्रिविम रसायन की व्याख्या कीजिए।
  \item फिर जल में सिल्वर नाइट्रेट से आयोडाइड हटाए जाते हैं। क्या बनता है?
  \item KCl मिलाने से सिसप्लैटिन मिलता है। विन्यास क्यों नहीं बदलता?
  \item इसके स्थान पर आप ट्रांसप्लैटिन कैसे बनाएँगे?
  \item ट्रांसप्लैटिन औषध के रूप में किसी काम का क्यों नहीं है, यद्यपि वह भी अभिक्रिया करता है?
\end{enumerate}

\medskip\noindent\textbf{भाग III --- ऐक्वाकरण।}
\begin{enumerate}[resume]
  \item प्रथम ऐक्वाकरण की अर्ध-आयु परिकलित कीजिए।
  \item क्लोराइड-मुक्त विलयन में \qty{2.0}{h} के बाद ऐक्वाकृत अंश परिकलित कीजिए।
  \item DNA के प्रति ऐक्वा संकुल क्रियाशील रूप क्यों है?
  \item ऐक्वाकरण वियोजी है या सहयोजी? उसका \omterm{def:b3:complex-mechanisms:activation-volume}{सक्रियण
        आयतन} क्या होगा?
  \item इंजेक्शन के लिए औषध विलयन लवण-जल में क्यों बनाया जाता है?
\end{enumerate}

\medskip\noindent\textbf{भाग IV --- प्लाज़्मा और कोशिका।}
\begin{enumerate}[resume]
  \item ऐक्वा संकुल का साम्य अंश
        $[\ce{Cl-}]$ के फलन के रूप में लिखिए।
  \item इसे प्लाज़्मा में परिकलित कीजिए।
  \item इसे कोशिका के भीतर परिकलित कीजिए।
  \item अनुपात परिकलित कीजिए।
  \item इसीलिए औषध मुख्यतः कोशिकाओं के भीतर ही क्यों कार्य करती है?
  \item कोशिका के भीतर कौन-से अन्य लिगैंड प्लैटिनम से बँधकर औषध को
        निष्क्रिय कर सकते हैं?
  \item दोनों ऐम्मीन लिगैंड पूरे समय बँधे क्यों रहते हैं?
  \item परिणाम बताइए: कोशिका के भीतर ऐक्वा-संकुल अंश का
        प्लाज़्मा वाले अंश से अनुपात।
\end{enumerate}
\end{problem}
